\documentclass[11pt]{scrartcl}

\usepackage[stixtwo]{fontsetup}
\usepackage{amsmath}
\usepackage{tikz}

\begin{document}
Az $1$ pont körüli forgatás: eltoljuk a síkot úgy, hogy az $1$ az origóba
kerüljön, az origó körül elforgatunk ($e^{i\varphi}$-vel szorzunk), majd
visszatoljuk. A $60^\circ$-os pozitív irányú forgatás az
$e^{i \frac{\pi}{3}}$-mal való szorzás, így
\[
	g(z)
	= 1 + e^{i \frac{\pi}{3}} (z - 1)
	= 1 + \left( \frac{1}{2} + \frac{\sqrt{3}}{2} i \right) (z - 1).
\]

Például a $z = 0$ pont képe
$g(0) = 1 - e^{i \frac{\pi}{3}} = \frac{1}{2} - \frac{\sqrt{3}}{2} i$:
\begin{center}
	\begin{tikzpicture}[scale=1.6, >=stealth]
		\draw[->] (-1, 0) -- (2, 0) node[right] {$\operatorname{Re}$};
		\draw[->] (0, -1.4) -- (0, 1) node[above] {$\operatorname{Im}$};
		\draw[dashed] (1, 0) circle (1);
		\draw (1, 0) -- (0, 0);
		\draw (1, 0) -- (0.5, -0.866);
		\draw[->] (0.6, 0) arc[start angle=180, end angle=240, radius=0.4];
		\node at (0.45, -0.3) {$60^\circ$};
		\fill (1, 0) circle (2pt) node[above] {$1$};
		\fill (0, 0) circle (2pt) node[above left] {$z = 0$};
		\fill (0.5, -0.866) circle (2pt) node[below] {$g(0)$};
	\end{tikzpicture}
\end{center}
\end{document}
